% System architecture figure — compile standalone or \input into main.tex
\begin{figure}[htbp]
\centering
\begin{tikzpicture}[
  node distance = 1.4cm and 2.2cm,
  box/.style  = {rectangle, rounded corners=5pt, draw=#1, fill=#1!10,
                 minimum width=2.8cm, minimum height=0.9cm,
                 font=\small\sffamily, align=center, thick},
  arr/.style  = {-{Stealth[length=6pt]}, thick, #1},
  lbl/.style  = {font=\footnotesize\sffamily\itshape},
]

% ── Left: Ingestion path ──────────────────────────────────────────────────────
\node[box=argblue, minimum width=3.6cm] (sources)
  {GSAS-II docs\\{\tiny 152 HTML (+179 book, opt.)}};

\node[box=argblue, below=of sources] (fetch)
  {Fetch \&\\Parse};

\node[box=argblue, below=of fetch] (chunk)
  {Chunk\\($\le$1200 chars)};

\node[box=argblue, below=of chunk] (embed)
  {Embed\\{\tiny bge-base-en-v1.5}};

\node[box=gray,   below=of embed,  minimum width=3.6cm] (chroma)
  {ChromaDB\\{\tiny 5{,}838 vectors}};

% ── Right: Query path ─────────────────────────────────────────────────────────
\node[box=apsorange, right=3.8cm of sources, minimum width=3.2cm] (user)
  {User question};

\node[box=apsorange, below=of user] (qembed)
  {Embed query\\{\tiny bge-base-en-v1.5}};

\node[box=apsorange, below=of qembed] (retrieve)
  {Retrieve top-6\\{\tiny cosine sim + URL dedup}};

\node[box=apsorange, below=of retrieve] (llm)
  {LLM generation\\{\tiny Ollama / llama-cpp-python}};

\node[box=apsorange, below=of llm, minimum width=3.2cm] (answer)
  {Cited answer\\{\tiny + source links [1..6]}};

% ── Arrows: ingestion ─────────────────────────────────────────────────────────
\draw[arr=argblue] (sources) -- (fetch);
\draw[arr=argblue] (fetch)   -- (chunk);
\draw[arr=argblue] (chunk)   -- (embed);
\draw[arr=argblue] (embed)   -- (chroma);

% ── Arrows: query ─────────────────────────────────────────────────────────────
\draw[arr=apsorange] (user)     -- (qembed);
\draw[arr=apsorange] (qembed)   -- (retrieve);
\draw[arr=apsorange] (retrieve) -- (llm);
\draw[arr=apsorange] (llm)      -- (answer);

% ── Cross arrow: DB → retrieve ────────────────────────────────────────────────
\draw[arr=black!60, dashed]
  (chroma.east) -- ++(0.6,0) |- (retrieve.west);

% ── Labels ────────────────────────────────────────────────────────────────────
\node[lbl, above=0.15cm of sources] {Ingestion (one-time)};
\node[lbl, above=0.15cm of user]    {Query (real-time)};

\end{tikzpicture}
\caption{\textit{Query-GSAS} system architecture.
  \textit{Left (blue):} The one-time ingestion pipeline fetches, parses,
  chunks, and embeds the GSAS-II documentation corpus into a local
  ChromaDB vector store using the \texttt{bge-base-en-v1.5} ONNX model.
  \textit{Right (orange):} At query time, the user's question is embedded
  with the same model; the top-30 candidates are retrieved by cosine
  similarity, deduplicated to one chunk per source URL, and the six
  most relevant are forwarded to an LLM for cited answer synthesis.
  All computation runs on the local machine.}
\label{fig:architecture}
\end{figure}
